\vspace{10pt}

\subsection{Sustainability}

The technological revolution has accelerated dramatically over the last few decades. A clear example is the launch of the first telephone (1876), the first computer (1937), the first laptop (1974), the first smartphone (2007), and the first artificial intelligence (2017) \cite{history-of-technology-timeline}. This trend reveals that the time it takes for a technological invention to emerge and significantly impact society continues to shrink. This exponential growth underscores the increasing importance of evaluating the economic, environmental, and social impacts of technological projects. These three dimensions together constitute sustainability, the topic discussed in this section.

Technology is shaping our future, at least in the coming years. The primary mission of every engineer is to make people's daily lives easier. However, as mentioned earlier, developing a solution is not enough; a study of its sustainability effects is also a must. Therefore, addressing sustainability questions enables a better evaluation of the project's impact across the economy, environment, and society. This analysis raises awareness of the weight of sustainability in project development.

The author's sustainability consciousness is quite limited. The knowledge extends only to well-known aspects such as environmental footprint and social impact. The economic dimension and multiple sustainability indicators require further study. The subsequent subsection answers a series of key sustainability questions within the author's current expertise.

\vspace{5pt}
\newpage

\subsubsection{Economic Impact}

\vspace{5pt}

\textbf{Regarding PPP: Have you estimated the project execution cost?}

The project budget is detailed in Section \ref{sec:budget}. All required costs have been considered, including personnel, hardware, software, contingencies, and incidents. Only indispensable resources are used and are included in the budget. 

The computer used throughout the project was recently purchased and remains under warranty. Since there is an overabundance of hardware resources in the office, they can be easily replaced if needed. Regarding the software resources, they were already available to the company and are virtual; as a consequence, fixes and repairs are not included in the budget. 

Additionally, only essential tasks demand a higher time cost. Since the AI model is connected to an external resource, this development phase can be significantly reduced while obtaining accurate data extraction. From an objective standpoint, a meticulous budget estimate has been established.

\vspace{5pt}

\textbf{Regarding Useful Life: How are currently solved economic issues (costs...) related to the problem that you want to address (state of the art)?}

The current expense management process imposes an excessive time cost for employees, who manually enter expenses into Excel spreadsheets. The same challenge is encountered in the finance department, which must verify the accuracy of each line item before approving reimbursements. Additionally, the current process is prone to human error regarding manual expense data entry and review. 

In short, this time-consuming, low-efficiency process delivers minimal value while misallocating company personnel costs for both employees and finance staff.

\vspace{5pt}

\textbf{Regarding Useful Life: How will your solution improve economic issues (costs, etc.) with respect to other existing solutions?}

The proposed solution automates the current cumbersome process. It provides automatic receipt data extraction and a more intuitive dashboard for displaying receipt information. This eliminates the tedious manual data entry, enhances visualization beyond traditional spreadsheets, and reduces human error. As a result, the workflow is significantly streamlined, delivering accurate receipt processing while maximizing efficiency. This enables employees to focus on higher-value tasks, thus maximizing personnel investment returns.

\subsubsection{Environmental Impact}

\textbf{Regarding PPP: Have you estimated the environmental impact of the project? }
The only resource directly affecting the environmental footprint is energy consumption. This resource is required throughout the entire project, as specified in Section 6.1.2. Both hardware resources and the shared workspace consume electricity continuously. However, sharing the workspace's lighting with other employees reduces its environmental impact.

\textbf{Regarding PPP: Did you plan to minimize its impact, for example, by reusing resources? }

The application uses Azure for software resources. Currently, the Azure resources required for application development cost less than traditional physical resources. This cost optimization results from significantly reduced resource consumption, as cloud resources eliminate the need for transportation and frequent physical replacements. Updates and fixes are performed virtually, further reducing hardware requirements and electronic waste.

\textbf{Regarding Useful Life: How is the problem that you want to address currently solved (state of the art)?}

The primary strategy for reducing resource usage is to develop the most efficient application possible, avoiding unnecessary resources and inefficient code that could overuse resources. This ensures that only essential resources are consumed each time employees use the application, limiting overall resource consumption throughout its lifecycle.

\textbf{Regarding Useful Life: How will your solution improve the environment with respect to other existing solutions? }

This application is tailored for the company's exclusive use and excludes any irrelevant functionality from the final product. It streamlines resource consumption by avoiding the unnecessary computation performed by other existing applications.



\subsubsection{Social Impact}

\vspace{5pt}

\textbf{Regarding PPP: What do you think you will achieve—in terms of personal growth—from doing this project?  }

This project enables the author to apply prior academic knowledge and software development skills in a real-world context. As a result, it enriches their working experience as future software engineer. Furthermore, it is the author's first experience managing and leading an entire project independently, strengthening self-learning, self-management skills and the ability to build software projects from scratch.

\vspace{5pt}

\textbf{Regarding Useful Life: How is the problem that you want to address currently solved (state of the art)?}

\newpage

The current expense management process is handled manually through traditional methods. This involves manually entering data into Excel spreadsheets, emailing them to the finance department, and reviewing them line by line. 

\vspace{5pt}

\textbf{Regarding Useful Life: How will your solution improve the quality of life (social dimension) with respect to other existing solutions? }

The proposed solution optimizes time spent on expense management while reducing employees' frustration and fatigue. This frees employees from this tedious task, allowing additional time to focus on more meaningful work. As a result, it maximizes employees' productivity.

\vspace{5pt}

\textbf{Regarding Useful Life: Is there a real need for the project?}

Regarding all project sustainability impacts (social, environmental, and economic), there is a clear need to build a robust solution. This solution overcomes challenges from all involved parties while introducing a sustainability-focused approach.
 


